\subsection{\texorpdfstring{Computation of the mappings \(\operatorname{Tor}^{\mathbf{A}}(g, f)\) and \(\operatorname{Ext}_{\mathbf{A}}(f, h)\)}{Computation of the mappings \textbackslash operatorname\{Tor\}\^{}\{\textbackslash mathbf\{A\}\}(g, f) and \textbackslash operatorname\{Ext\}\_\{\textbackslash mathbf\{A\}\}(f, h)}}

Let \(f: \mathbf{M} \to \mathbf{M}'\) , \(h: \mathbf{N}' \to \mathbf{N}\) be homomorphisms of left A-modules, \(g: \mathbf{P} \to \mathbf{P}'\) a homomorphism of right A-modules, \(a: \mathbf{R} \to \mathbf{M}\) , \(a': \mathbf{R}' \to \mathbf{M}'\) , \(b: \mathbf{S} \to \mathbf{P}\) , \(b': \mathbf{S}' \to \mathbf{P}'\) left resolutions of M, M', P, P', respectively, \(c: \mathbf{N} \to \mathbf{E}\) , \(c': \mathbf{N}' \to \mathbf{E}'\) right resolutions of N and N', \(\tilde{f}: \mathbf{R} \to \mathbf{R}'\) , \(\tilde{g}: \mathbf{S} \to \mathbf{S}'\) , \(\tilde{h}: \mathbf{E}' \to \mathbf{E}\) morphisms of complexes such that

\[
a ^ {\prime} \circ \tilde {f} = f \circ a, b ^ {\prime} \circ \tilde {g} = g \circ b, \tilde {h} \circ c ^ {\prime} = c \circ h.
\]

PROPOSITION 2. --- The following two diagrams are commutative:

\[
\begin{array}{c} \operatorname{Tor} ^ {\mathbf {A}} (\mathrm{P}, \mathrm{M}) \xrightarrow {\psi (\mathrm{S} , \mathrm{R})} \mathrm{H} (\mathrm{S} \otimes_ {\mathbf {A}} \mathrm{R}) \\ \operatorname{Tor} ^ {\mathbf {A}} (g, f) \Big \downarrow \qquad \qquad \qquad \qquad \qquad \qquad \qquad \qquad \qquad \qquad \qquad \qquad \qquad \qquad \qquad \qquad \qquad \qquad \qquad \qquad \qquad \qquad \qquad \qquad \qquad \qquad \qquad \qquad \qquad \qquad \qquad \qquad \qquad \qquad \\ \operatorname{Tor} ^ {\mathbf {A}} (\mathrm{P} ^ {\prime}, \mathrm{M} ^ {\prime}) \xrightarrow {\psi (\mathrm{S} ^ {\prime} , \mathrm{R} ^ {\prime})} \mathrm{H} (\mathrm{S} ^ {\prime} \otimes_ {\mathbf {A}} \mathrm{R} ^ {\prime})  , \end{array}
\]

\[
\begin{array}{c} \mathrm{H} (\operatorname{Homgr} _ {\mathbf {A}} (\mathrm{R} ^ {\prime}, \mathrm{E} ^ {\prime})) \xrightarrow {\varphi (\mathrm{R} ^ {\prime} , \mathrm{E} ^ {\prime})} \operatorname{Ext} _ {\mathbf {A}} (\mathrm{M} ^ {\prime}, \mathrm{N} ^ {\prime}) \\ \mathrm{H} (\operatorname{Homgr} _ {\mathbf {A}} (\tilde {f}, \tilde {h})) \Biggl \downarrow \\ \mathrm{H} (\operatorname{Homgr} _ {\mathbf {A}} (\mathrm{R}, \mathrm{E})) \xrightarrow {\varphi (\mathrm{R} , \mathrm{E})} \operatorname{Ext} _ {\mathbf {A}} (\mathrm{M}, \mathrm{N}). \end{array}
\]

Let \(\alpha: L(M) \to R\) , \(\alpha': L(M') \to R'\) , \(\gamma: L(P) \to S\) , \(\gamma': L(P') \to S'\) morphisms of complexes such that

\[
a \circ \alpha = p _ {\mathrm{M}}, a ^ {\prime} \circ \alpha^ {\prime} = p _ {\mathrm {M^ {\prime}}}, b \circ \gamma = p _ {\mathrm{P}}, b ^ {\prime} \circ \gamma^ {\prime} = p _ {\mathrm {P^ {\prime}}}.
\]

By definition, \(\mathrm{H}(\tilde{g}\otimes\tilde{f})\circ\psi(\mathrm{S},\mathrm{R})\) is equal to

\[
\mathrm{H} (\tilde {g} \otimes \tilde {f}) \circ \mathrm{H} (\gamma \otimes \alpha) = \mathrm{H} ((\tilde {g} \circ \gamma) \otimes (\tilde {f} \circ \alpha)),
\]

while \(\psi(\mathbf{S}',\mathbf{R}')\circ\mathrm{Tor}^{\mathbf{A}}(g,f)\) is equal to

\[
\mathrm{H} \left(\gamma^ {\prime} \otimes \alpha^ {\prime}\right) \circ \mathrm{H} \left(\mathrm{L} (g) \otimes \mathrm{L} (f)\right) = \mathrm{H} \left(\left(\gamma^ {\prime} \circ \mathrm{L} (g)\right) \otimes \left(\alpha^ {\prime} \circ \mathrm{L} (f)\right)\right).
\]

On the other hand, \(\alpha' \circ L(f)\) and \(\tilde{f} \circ \alpha\) are two morphisms from L(M) into R′ such that \(a' \circ (\alpha' \circ L(f)) = p_{\mathrm{M}} \circ L(f) = f \circ p_{\mathrm{M}} = f \circ a \circ \alpha = a' \circ (\tilde{f} \circ \alpha)\) . By X, p.~49, prop. 3, \(\alpha' \circ L(f)\) and \(\tilde{f} \circ \alpha\) are homotopic; likewise \(\gamma' \circ L(g)\) and \(\tilde{g} \circ \gamma\) are homotopic, and consequently so are \((\gamma' \circ L(g)) \otimes (\alpha' \circ L(f))\) and \((\tilde{g} \circ \gamma) \otimes (\tilde{f} \circ \alpha)\) by prop. 3 of X, p.~64. We therefore have

\[
\begin{array}{r l} \mathrm{H} (\tilde {g} \otimes \tilde {f}) \circ \psi (\mathrm{S}, \mathrm{R}) & = \mathrm{H} ((\tilde {g} \circ \gamma) \otimes (\tilde {f} \circ \alpha)) = \mathrm{H} ((\gamma^ {\prime} \circ \mathrm{L} (g)) \otimes (\alpha^ {\prime} \circ \mathrm{L} (f))) \\ & = \psi (\mathrm{S} ^ {\prime}, \mathrm{R} ^ {\prime}) \circ \mathrm{Tor} ^ {\mathrm{A}} (g, f). \end{array}
\]

We argue analogously for the second diagram.

Remark. --- Consider similarly commutative diagrams of morphisms of complexes.

\[
\begin{array}{c c c} \mathrm{R} _ {1} \xrightarrow {a _ {1}} \mathrm{M} & \mathrm{S} _ {1} \xrightarrow {b _ {1}} \mathrm{P} & \mathrm{N} ^ {\prime} \xrightarrow {c _ {1} ^ {\prime}} \mathrm{E} ^ {\prime} \\ \tilde {f} \Big \downarrow & f \Big \downarrow & g \Big \downarrow \\ \mathrm{R} _ {1} ^ {\prime} \xrightarrow {a _ {1} ^ {\prime}} \mathrm{M} ^ {\prime} & \mathrm{S} _ {1} ^ {\prime} \xrightarrow {b _ {1} ^ {\prime}} \mathrm{P} ^ {\prime} & \mathrm{N} \xrightarrow {c _ {1}} \mathrm{E} _ {1} \end{array}
\]

where the complexes \(R_{1}\) , \(R_{1}^{\prime}\) , \(S_{1}\) , \(S_{1}^{\prime}\) are projective and zero on the right, and where the complexes \(E_{1}\) , \(E_{1}^{\prime}\) are injective and zero on the left. Then

\[
\begin{array}{r l} & {\mathrm{Tor} ^ {\mathbf {A}} (g, f) \circ \psi^ {\prime} (\mathrm{S} _ {1}, \mathrm{R} _ {1}) = \psi^ {\prime} (\mathrm{S} _ {1} ^ {\prime}, \mathrm{R} _ {1} ^ {\prime}) \circ \mathrm{H} (\tilde {g} \otimes \tilde {f})} \\ & {\varphi^ {\prime} (\mathrm{R} _ {1}, \mathrm{E} _ {1}) \circ \mathrm{Ext} _ {\mathbf {A}} (f, h) = \mathrm{H} (\mathrm{Homgr} _ {\mathbf {A}} (\tilde {f}, \tilde {h})) \circ \varphi^ {\prime} (\mathrm{R} _ {1} ^ {\prime}, \mathrm{E} _ {1} ^ {\prime})} \end{array}
\]

as is shown by an argument analogous to that of Prop. 2.

